\chapter{مقدمه‌ای بر داده‌کاوی}
\label{ch:intro}

\section*{اهداف این فصل}
\begin{itemize}
    \item درک تفاوت داده‌کاوی با آمار سنتی، پایگاه داده و یادگیری ماشین
    \item آشنایی با فرآیند کشف دانش از داده و جایگاه داده‌کاوی در آن
    \item شناخت انواع الگوهایی که می‌توان از داده استخراج کرد
    \item آشنایی با کاربردهای واقعی و چالش‌های اصلی داده‌کاوی
\end{itemize}

\section{داده‌کاوی چیست؟}

در دهه‌های اخیر، حجم داده‌های تولیدشده توسط سازمان‌ها، وب‌سایت‌ها، شبکه‌های اجتماعی،
حسگرها و سامانه‌های تراکنشی به شکلی بی‌سابقه رشد کرده است. بانک‌ها میلیون‌ها تراکنش
روزانه ثبت می‌کنند، فروشگاه‌های آنلاین رفتار خرید هر کاربر را ذخیره می‌کنند و شبکه‌های
اجتماعی روزانه میلیاردها پیام و تعامل تولید می‌کنند. مسئلهٔ اصلی این نیست که داده کم
است؛ برعکس، مسئله این است که در میان این حجم عظیم داده، دانش مفید و قابل‌اتکا چگونه
استخراج شود. این دقیقاً همان مسئله‌ای است که \dmterm{\textbf{داده‌کاوی}}{Data Mining}
به آن می‌پردازد.

داده‌کاوی را می‌توان این‌گونه تعریف کرد:

\begin{quote}
داده‌کاوی فرآیند کشف الگوها، روابط و دانش پنهان، غیربدیهی و بالقوه مفید از میان
حجم زیادی از داده است.
\end{quote}

سه ویژگی در این تعریف کلیدی هستند:

\begin{itemize}
    \item \textbf{پنهان بودن}: الگوهایی که به‌راحتی و با نگاه ساده به داده قابل مشاهده نیستند.
    \item \textbf{غیربدیهی بودن}: نتیجه نباید چیزی باشد که از قبل بدیهی یا شناخته‌شده است
          (مثلاً «مشتریانی که نان می‌خرند، معمولاً نان می‌خواهند» یک الگوی بی‌ارزش است).
    \item \textbf{مفید بودن}: الگوی کشف‌شده باید قابل استفاده در تصمیم‌گیری یا پیش‌بینی باشد.
\end{itemize}

\subsection{داده‌کاوی در برابر آمار، پایگاه داده و یادگیری ماشین}

سؤالی که معمولاً برای دانشجو پیش می‌آید این است: داده‌کاوی چه تفاوتی با آماری که
پیش از این خوانده‌اید، یا با پایگاه داده و یادگیری ماشین دارد؟ پاسخ این است که
داده‌کاوی یک حوزهٔ میان‌رشته‌ای است که از هر سه این حوزه‌ها بهره می‌برد:

\begin{itemize}
    \item از \textbf{آمار}، مبانی نظری برای اندازه‌گیری عدم قطعیت، آزمون فرضیه و
          برازش مدل را وام می‌گیرد.
    \item از \textbf{سیستم‌های پایگاه داده}، فناوری لازم برای ذخیره، بازیابی و
          مدیریت کارآمد حجم بزرگی از داده را می‌گیرد که اغلب در حافظه جا نمی‌شود.
    \item از \textbf{یادگیری ماشین} و هوش مصنوعی، الگوریتم‌هایی برای یادگیری خودکار
          الگو از داده به‌جای برنامه‌نویسی دستی قوانین را به کار می‌گیرد.
\end{itemize}

تفاوت اصلی داده‌کاوی با یادگیری ماشین صرف، در تمرکز آن است: یادگیری ماشین بیشتر
به الگوریتم و بهینه‌سازی مدل می‌پردازد، در حالی که داده‌کاوی کل فرآیند — از
داده خام تا دانش قابل استفاده — را در نظر می‌گیرد. جدول \ref{tab:dm-comparison}
این مقایسه را به‌طور خلاصه نشان می‌دهد.

\begin{table}[htbp]
\centering
\caption{مقایسهٔ داده‌کاوی با حوزه‌های نزدیک به آن}
\label{tab:dm-comparison}
\begin{tabular}{@{}p{2.6cm}p{5.8cm}p{4.6cm}@{}}
\toprule
\textbf{حوزه} & \textbf{تمرکز اصلی} & \textbf{سهم آن در داده‌کاوی} \\
\midrule
آمار & اندازه‌گیری عدم قطعیت، آزمون فرضیه، برازش مدل & مبانی نظری استنباط از داده \\
\addlinespace
پایگاه داده & ذخیره، بازیابی و مدیریت کارآمد داده حجیم & زیرساخت دسترسی به داده \\
\addlinespace
یادگیری ماشین & الگوریتم‌های یادگیری خودکار الگو از داده & موتور کشف الگو \\
\addlinespace
داده‌کاوی & کل فرآیند: از داده خام تا دانش قابل استفاده & ترکیب هر سه حوزهٔ بالا \\
\bottomrule
\end{tabular}
\end{table}

\section{فرآیند کشف دانش از داده}
\label{sec:kdd}

داده‌کاوی معمولاً به‌اشتباه به‌عنوان مترادف کل فرآیند تحلیل داده به کار می‌رود،
در حالی که در واقع تنها یکی از مراحل یک فرآیند بزرگ‌تر به نام
\dmterm{\textbf{کشف دانش از داده}}{Knowledge Discovery in Databases (KDD)} است.
شکل \ref{fig:kdd-flow} این فرآیند هفت‌مرحله‌ای را نشان می‌دهد.

\begin{figure}[htbp]
\centering
\begin{tikzpicture}[
    node distance=0.55cm and 0cm,
    box/.style={rectangle, rounded corners, draw=blue!60, thick,
                fill=blue!8, minimum width=8.5cm, minimum height=0.9cm,
                align=center, font=\small},
    arr/.style={-{Latex[length=2mm]}, thick, draw=blue!60}
]
\node[box] (n1) {گام ۱: پاک‌سازی داده};
\node[box, below=of n1] (n2) {گام ۲: یکپارچه‌سازی داده};
\node[box, below=of n2] (n3) {گام ۳: انتخاب داده};
\node[box, below=of n3] (n4) {گام ۴: تبدیل داده};
\node[box, below=of n4, fill=orange!15, draw=orange!70] (n5) {گام ۵: داده‌کاوی};
\node[box, below=of n5] (n6) {گام ۶: ارزیابی الگو};
\node[box, below=of n6, fill=green!10, draw=green!60!black] (n7) {گام ۷: ارائهٔ دانش};

\draw[arr] (n1) -- (n2);
\draw[arr] (n2) -- (n3);
\draw[arr] (n3) -- (n4);
\draw[arr] (n4) -- (n5);
\draw[arr] (n5) -- (n6);
\draw[arr] (n6) -- (n7);
\end{tikzpicture}
\caption{مراحل هفت‌گانهٔ فرآیند کشف دانش از داده}
\label{fig:kdd-flow}
\end{figure}

جزئیات هر مرحله به شرح زیر است:

\begin{enumerate}
    \item \dmterm{\textbf{پاک‌سازی داده}}{Data Cleaning}: حذف نویز و داده‌های ناسازگار.
    \item \dmterm{\textbf{یکپارچه‌سازی داده}}{Data Integration}: ترکیب داده از چند منبع مختلف.
    \item \dmterm{\textbf{انتخاب داده}}{Data Selection}: بازیابی داده‌های مرتبط با هدف تحلیل.
    \item \dmterm{\textbf{تبدیل داده}}{Data Transformation}: تبدیل داده به قالبی مناسب برای کاوش
          (مثلاً از طریق تجمیع یا نرمال‌سازی).
    \item \textbf{داده‌کاوی}: به‌کارگیری روش‌های هوشمند برای استخراج الگو.
    \item \dmterm{\textbf{ارزیابی الگو}}{Pattern Evaluation}: شناسایی الگوهای واقعاً جالب بر اساس
          معیارهایی مانند \dmterm{جالب بودن}{Interestingness}.
    \item \dmterm{\textbf{ارائهٔ دانش}}{Knowledge Presentation}: نمایش دانش کشف‌شده به کاربر با
          استفاده از تکنیک‌های تجسم داده و گزارش‌گیری.
\end{enumerate}

نکتهٔ مهم این است که در عمل، این مراحل خطی نیستند؛ فرآیند کشف دانش تکرارشونده است و
معمولاً پس از ارزیابی نتایج اولیه، لازم است به مراحل پیش‌تر (مثلاً انتخاب داده
یا تبدیل داده) بازگشت. در بسیاری از منابع صنعتی، نسخهٔ ساده‌شده‌تری از این
چرخه با نام \dmterm{\LR{CRISP-DM}}{Cross-Industry Standard Process for Data Mining}
نیز رایج است که شش مرحله دارد و به‌جای یک خط مستقیم، به‌صورت یک \textbf{چرخه}
دیده می‌شود (شکل \ref{fig:crispdm-cycle})، زیرا معمولاً پس از استقرار یک مدل در
عمل، پروژه دوباره از مرحلهٔ درک کسب‌وکار آغاز می‌شود.

\begin{figure}[htbp]
\centering
\begin{tikzpicture}[
    every node/.style={font=\scriptsize},
    circnode/.style={rectangle, rounded corners, draw=purple!70, thick,
                      fill=purple!8, minimum width=2.7cm, minimum height=0.85cm, align=center}
]
\def\n{6}
\def\R{3.3cm}
\node[circnode] (c0) at (90:\R)    {درک کسب‌وکار};
\node[circnode] (c1) at (30:\R)    {درک داده};
\node[circnode] (c2) at (-30:\R)   {آماده‌سازی داده};
\node[circnode] (c3) at (-90:\R)   {مدل‌سازی};
\node[circnode] (c4) at (-150:\R)  {ارزیابی};
\node[circnode] (c5) at (150:\R)   {استقرار};

\draw[-{Latex[length=2mm]}, thick, purple!70] (c0) -- (c1);
\draw[-{Latex[length=2mm]}, thick, purple!70] (c1) -- (c2);
\draw[-{Latex[length=2mm]}, thick, purple!70] (c2) -- (c3);
\draw[-{Latex[length=2mm]}, thick, purple!70] (c3) -- (c4);
\draw[-{Latex[length=2mm]}, thick, purple!70] (c4) -- (c5);
\draw[-{Latex[length=2mm]}, thick, purple!70] (c5) -- (c0);
\end{tikzpicture}
\caption{چرخهٔ \LR{CRISP-DM} به‌عنوان نسخهٔ صنعتی و تکرارشوندهٔ فرآیند کشف دانش}
\label{fig:crispdm-cycle}
\end{figure}

\subsection{چرا پیش‌پردازش این‌قدر مهم است؟}

تجربهٔ عملی نشان می‌دهد که در اغلب پروژه‌های داده‌کاوی، بیش از ۶۰ تا ۷۰ درصد
زمان صرف مراحل پاک‌سازی، یکپارچه‌سازی و تبدیل داده می‌شود، نه اجرای الگوریتم
یادگیری. دلیل آن ساده است: کیفیت هر مدل داده‌کاوی مستقیماً به کیفیت داده‌ای
که به آن داده می‌شود بستگی دارد؛ اصلی که با عبارت معروف «زباله وارد شود،
زباله بیرون می‌آید»\footnote{\LR{Garbage In, Garbage Out}} شناخته می‌شود. جزئیات
این مرحله در فصل \ref{ch:data} بررسی خواهد شد.

\section{انواع الگوهای قابل کشف}

خروجی داده‌کاوی می‌تواند اشکال گوناگونی داشته باشد. به‌طور کلی، وظایف داده‌کاوی
به دو دستهٔ اصلی تقسیم می‌شوند:

\begin{itemize}
    \item \dmterm{\textbf{وظایف توصیفی}}{Descriptive}: هدف، شناخت ویژگی‌های کلی داده است،
          بدون آنکه لزوماً پیش‌بینی خاصی صورت گیرد. نمونه‌ها: خوشه‌بندی، قوانین
          انجمنی، تشخیص داده پرت.
    \item \dmterm{\textbf{وظایف پیش‌بینی‌کننده}}{Predictive}: هدف، استفاده از برخی
          ویژگی‌ها برای پیش‌بینی مقدار یک ویژگی دیگر است. نمونه‌ها: طبقه‌بندی،
          رگرسیون.
\end{itemize}

جدول \ref{tab:dm-tasks} خلاصه‌ای از وظایف اصلی داده‌کاوی و فصل مرتبط با هر یک
در این کتاب را نشان می‌دهد.

\begin{table}[htbp]
\centering
\caption{وظایف اصلی داده‌کاوی و فصل مرتبط با هر یک در این کتاب}
\label{tab:dm-tasks}
\begin{tabular}{@{}lll@{}}
\toprule
\textbf{وظیفه} & \textbf{نوع} & \textbf{فصل مرتبط} \\
\midrule
طبقه‌بندی (\LR{Classification}) & پیش‌بینی‌کننده & فصل \ref{ch:classification-basic}، \ref{ch:classification-advanced} \\
رگرسیون (\LR{Regression}) & پیش‌بینی‌کننده & فصل \ref{ch:regression} \\
قوانین انجمنی (\LR{Association Rules}) & توصیفی & فصل \ref{ch:association} \\
خوشه‌بندی (\LR{Clustering}) & توصیفی & فصل \ref{ch:clustering} \\
تشخیص داده پرت (\LR{Outlier Detection}) & توصیفی/پیش‌بینی‌کننده & فصل \ref{ch:outliers} \\
\bottomrule
\end{tabular}
\end{table}

\section{کاربردهای داده‌کاوی}

داده‌کاوی امروزه در تقریباً هر صنعتی کاربرد دارد. جدول \ref{tab:dm-applications}
چند نمونهٔ رایج آن را نشان می‌دهد.

\begin{table}[htbp]
\centering
\caption{نمونه‌هایی از کاربرد داده‌کاوی در صنایع مختلف}
\label{tab:dm-applications}
\begin{tabular}{@{}p{3.2cm}p{9.8cm}@{}}
\toprule
\textbf{حوزه} & \textbf{کاربرد نمونه} \\
\midrule
بانکداری و مالی & تشخیص تقلب در تراکنش‌های کارت اعتباری، \dmterm{امتیازدهی اعتباری}{Credit Scoring} مشتریان \\
\addlinespace
خرده‌فروشی & \dmterm{تحلیل سبد خرید}{Market Basket Analysis} برای چیدمان محصولات و پیشنهاد خرید \\
\addlinespace
مراقبت سلامت & پیش‌بینی ریسک بیماری بر اساس سوابق بیمار، کشف الگوهای درمانی مؤثر \\
\addlinespace
بازاریابی & \dmterm{بخش‌بندی مشتریان}{Customer Segmentation}، \dmterm{سامانه‌های توصیه‌گر}{Recommender Systems} \\
\addlinespace
امنیت & \dmterm{تشخیص نفوذ به شبکه}{Intrusion Detection} \\
\bottomrule
\end{tabular}
\end{table}

\section{چالش‌های داده‌کاوی}

با وجود پیشرفت‌های چشمگیر، داده‌کاوی همچنان با چالش‌های جدی روبه‌روست:

\begin{itemize}
    \item \textbf{مقیاس‌پذیری}: بسیاری از الگوریتم‌های کلاسیک برای داده‌های
          کوچک طراحی شده‌اند و روی داده‌های بسیار بزرگ\footnote{\LR{Big Data}}
          کارایی کافی ندارند.
    \item \textbf{ابعاد بالا}: با افزایش تعداد ویژگی‌ها، بسیاری از روش‌ها
          دچار \dmterm{«نفرین ابعاد»}{Curse of Dimensionality} می‌شوند.
    \item \textbf{ناهمگونی و پیچیدگی داده}: داده‌های واقعی اغلب ترکیبی از
          متن، تصویر، سری زمانی و داده رابطه‌ای هستند.
    \item \textbf{حریم خصوصی و اخلاق داده}: استخراج الگو از داده‌های شخصی
          مسئولیت‌های حقوقی و اخلاقی جدی به همراه دارد.
    \item \textbf{تفسیرپذیری}: مدل‌های پیچیده (مانند جنگل تصادفی یا شبکه‌های
          عصبی) اغلب به‌صورت جعبه سیاه عمل می‌کنند و توضیح تصمیم آن‌ها دشوار است.
\end{itemize}

\section*{پیاده‌سازی در پایتون}

برای آشنایی عملی، مجموعه دادهٔ کلاسیک \LR{Iris} را بارگذاری کرده و یک
تحلیل اکتشافی دادهٔ ساده\footnote{\LR{Exploratory Data Analysis}} روی آن
انجام می‌دهیم: رسم پراکندگی
دو ویژگی «طول کاسبرگ» و «عرض کاسبرگ» برای سه گونهٔ گل، پیش از هرگونه
مدل‌سازی.

\begin{lstlisting}
import matplotlib.pyplot as plt
from sklearn.datasets import load_iris

iris = load_iris()
X, y, names = iris.data, iris.target, iris.target_names

fig, ax = plt.subplots(figsize=(5.5, 4.2))
for i, name in enumerate(names):
    mask = y == i
    ax.scatter(X[mask, 0], X[mask, 1], label=name.capitalize())

ax.set_xlabel("Sepal length (cm)")
ax.set_ylabel("Sepal width (cm)")
ax.legend()
plt.show()
\end{lstlisting}

خروجی این کد در شکل \ref{fig:iris-scatter} نشان داده شده است. حتی پیش از
اجرای هر الگوریتم داده‌کاوی، همین نمودار ساده به‌خوبی نشان می‌دهد که گونهٔ
\LR{Setosa} (نقاط آبی) بر اساس همین دو ویژگی به‌خوبی از دو گونهٔ دیگر
قابل تفکیک است، در حالی که \LR{Versicolor} و \LR{Virginica} همپوشانی
بیشتری دارند — دقیقاً نوع بینشی که تحلیل اکتشافی داده پیش از داده‌کاوی
باید فراهم کند.

\begin{figure}[htbp]
\centering
\includegraphics[width=0.72\linewidth]{figures/ch01_iris_exploratory.pdf}
\caption{نمودار پراکندگی طول و عرض کاسبرگ برای سه گونهٔ مجموعه دادهٔ Iris}
\label{fig:iris-scatter}
\end{figure}

کد کامل‌تر این فصل، به‌همراه چند نمودار اکتشافی دیگر، در بخش «فصل ۱» پیوست
الف (\ref{app:python}) آمده است.

\section*{خلاصهٔ فصل}

در این فصل با تعریف داده‌کاوی، جایگاه آن در فرآیند بزرگ‌تر کشف دانش، دسته‌بندی
وظایف اصلی آن به دو گروه توصیفی و پیش‌بینی‌کننده، و همچنین کاربردها و
چالش‌های اصلی این حوزه آشنا شدیم. در فصل بعد، به یکی از مهم‌ترین و
زمان‌برترین بخش‌های هر پروژهٔ داده‌کاوی، یعنی آشنایی با داده و پیش‌پردازش
آن، خواهیم پرداخت.

\section*{تمرین‌ها}

\begin{enumerate}
    \item تفاوت میان یک الگوی «بدیهی» و یک الگوی «جالب» را با یک مثال از
          حوزهٔ دلخواه خودتان توضیح دهید.
    \item برای هر یک از مراحل هفت‌گانهٔ کشف دانش (شکل \ref{fig:kdd-flow})، یک
          مثال عملی از یک پروژهٔ فرضی (مثلاً تحلیل رفتار مشتریان یک
          فروشگاه آنلاین) بیاورید.
    \item چرا چرخهٔ \LR{CRISP-DM} (شکل \ref{fig:crispdm-cycle}) به‌صورت
          دایره‌ای ترسیم می‌شود و نه به‌صورت یک خط مستقیم مانند فرآیند
          کشف دانش؟
    \item دو مورد از وظایف داده‌کاوی (مثلاً طبقه‌بندی و خوشه‌بندی) را از
          نظر نوع (توصیفی/پیش‌بینی‌کننده) و نوع داده‌ای که نیاز دارند
          (برچسب‌دار یا بدون برچسب) با هم مقایسه کنید.
    \item کد بخش «پیاده‌سازی در پایتون» را اجرا کنید و به‌جای دو ویژگی
          «طول و عرض کاسبرگ»، «طول و عرض گلبرگ» را رسم کنید. آیا تفکیک
          سه گونه در این حالت بهتر است یا بدتر؟ چرا؟
\end{enumerate}
